\chapter{Código para el nodo sensor}
\label{anexo:codigo_nodo}
A continuación se proporciona el enlace al repositorio que contiene el código fuente en C/C++ (Arduino) desarrollado para la placa Arduino UNO del nodo sensor. Este código incluye las rutinas de lectura de los sensores de CO\(_2\), temperatura, humedad, el control del actuador PWM y el empaquetado de la telemetría.

\begin{center}
    URL del repositorio: \href{https://github.com/sebasandre2020/Tesis-2026-1/tree/appCode/Codes/Nodes/main}{https://github.com/sebasandre2020/Tesis-2026-1/tree/appCode/Codes/Nodes/main}
\end{center}

\chapter{Código para el gateway personalizado}
\label{anexo:codigo_gateway}
Este anexo contiene el enlace al código fuente en C/C++ desarrollado para el microcontrolador ESP32 que actúa como gateway. Incluye las clases \texttt{LoRaManager} y \texttt{AzureIoTManager} encargadas de la recepción LoRa en modo no bloqueante y la publicación MQTT segura hacia Azure IoT Hub.

\begin{center}
    URL del repositorio: \href{https://github.com/sebasandre2020/Tesis-2026-1/tree/appCode/Codes/Gateway/main}{https://github.com/sebasandre2020/Tesis-2026-1/tree/appCode/Codes/Gateway/main}
\end{center}

\chapter{Código para la aplicación web}
\label{anexo:codigo_web}
A continuación se provee el enlace al código fuente de la plataforma web, subdividido en el \textit{backend} (desarrollado en .NET Core 8.0 utilizando Clean Architecture) y el \textit{frontend} (desarrollado en React).

\begin{center}
    URL del repositorio: \href{https://github.com/sebasandre2020/Tesis-2026-1/tree/appCode/Codes/Frontend}{https://github.com/sebasandre2020/Tesis-2026-1/tree/appCode/Codes/Frontend}
\end{center}

\chapter{Código IaC para infraestructura en Azure}
\label{anexo:codigo_iac}
Este anexo contiene el enlace a los \textit{scripts} y plantillas de Infraestructura como Código (IaC), utilizados para el aprovisionamiento y configuración automatizada de los recursos en Microsoft Azure (IoT Hub, SQL Database, Container Apps, Azure Functions, Static Web Apps, etc.).

\begin{center}
    URL del repositorio: \href{https://github.com/sebasandre2020/Tesis-2026-1/tree/appCode/Codes/Software\%20Infraestructure}{https://github.com/sebasandre2020/Tesis-2026-1/tree/appCode/Codes/Software\%20Infraestructure}
\end{center}

\chapter{Acceso a la aplicación web desplegada}
\label{anexo:acceso_web}
Este anexo proporciona el enlace de acceso directo a la aplicación web de monitoreo y alerta ambiental desplegada en producción en Microsoft Azure.

\begin{center}
    URL de la aplicación: \href{https://orange-pond-0bd85900f.7.azurestaticapps.net/}{https://orange-pond-0bd85900f.7.azurestaticapps.net/}
\end{center}